\chapter{Del AST al \emph{matcher}}
\label{cap:ast}

\begin{intuicion}
La pieza intermedia entre el parser y el matcher es la
\texttt{Hypothesis} ---el AST de la DSL, simplificado para que el
matcher pueda consumirlo sin parsear strings---. Este capítulo
recorre la struct \texttt{Hypothesis} campo por campo, muestra cómo
el matcher dispatcha según \texttt{shape} y \texttt{steps}, y
discute cómo otros transports (HTTP, MCP, programmatic) pueden
saltarse el parser y construir \texttt{Hypothesis} directamente.
\end{intuicion}

\section{La struct \texttt{Hypothesis}}

\begin{lstlisting}[language=Rust]
pub struct Hypothesis {
    pub name: String,
    pub steps: Vec<HypothesisStep>,
    pub k_simplicity: usize,
    pub aggregations: Vec<AggClause>,
    pub window_seconds: Option<i64>,
    pub shape: HypothesisShape,
}

pub enum HypothesisShape {
    Chain,
    Graph,
}

pub struct HypothesisStep {
    pub source_var: VarId,
    pub source_type: EntityType,
    pub source_predicates: Vec<Predicate>,
    pub rel_type: RelationType,
    pub edge_predicates: Vec<Predicate>,
    pub dest_var: VarId,
    pub dest_type: EntityType,
    pub dest_predicates: Vec<Predicate>,
}
\end{lstlisting}

\noindent
Notar las diferencias con la DSL textual:

\begin{itemize}
    \item Los \texttt{binding} (textuales: \texttt{u}, \texttt{h1})
    se convierten a \texttt{VarId} (numérico, único por
    hipótesis). El nombre se descarta; sólo importa la
    \emph{identidad}.
    \item Los predicados están parseados a \texttt{Predicate} (un
    enum); el matcher los evalúa sin re-parsear strings.
    \item \texttt{shape} es explícito: \texttt{Chain} o
    \texttt{Graph}, no derivado de la presencia de \texttt{;}.
    \item \texttt{aggregations} es lista plana de cláusulas
    HAVING; el evaluador las aplica al resultado.
\end{itemize}

\section{El \texttt{VarId}}

\begin{lstlisting}[language=Rust]
#[derive(Copy, Clone, PartialEq, Eq, Hash)]
pub struct VarId(pub usize);

impl Hypothesis {
    pub fn vars(&self) -> Vec<VarId> {
        let mut seen = HashSet::new();
        for step in &self.steps {
            seen.insert(step.source_var);
            seen.insert(step.dest_var);
        }
        seen.into_iter().collect()
    }
}
\end{lstlisting}

\noindent
\texttt{VarId} es un índice en el \emph{binding array} del matcher.
Para una chain de \(k\) pasos sin bindings compartidos: \(k + 1\)
\texttt{VarId}s distintos. Para una hipótesis graph donde el
binding \(c\) aparece en 3 cláusulas: \(c\) tiene un solo
\texttt{VarId}, los otros vértices tienen los suyos.

El parser asigna \texttt{VarId}s en orden de aparición, mapping
nombres textuales (\texttt{u}, \texttt{c}, \texttt{h1}) a índices
\([0, 1, 2, \dots]\). Los nombres textuales se mantienen \emph{sólo}
para el debugging y para el HAVING (que referencia bindings por
nombre).

\section{El campo \texttt{shape}}

\begin{lstlisting}[language=Rust]
pub enum HypothesisShape {
    Chain,
    Graph,
}
\end{lstlisting}

\noindent
La pieza nueva que destrabó L1 stage 3. El matcher branchea según
\texttt{shape} en \texttt{search\_temporal\_pattern\_yannakakis}:

\begin{lstlisting}[language=Rust]
match hypothesis.shape {
    HypothesisShape::Chain => {
        let bags = build_bags_chain(graph, hypothesis);
        semijoin_reduce_chain(&mut bags);
        enumerate_chain(&bags, hypothesis)
    }
    HypothesisShape::Graph => {
        let join_tree = JoinTree::build(hypothesis);
        let bags = build_bags_graph(graph, hypothesis, &join_tree);
        semijoin_reduce_tree(&mut bags, &join_tree);
        enumerate_tree(&bags, &join_tree, hypothesis)
    }
}
\end{lstlisting}

\noindent
Dos paths distintos pero \emph{paralelos}:
\begin{itemize}
    \item Chain (capítulo~\ref{cap:yannakakis-chain}): bags
    indexados por step (orden lineal); semi-join reducer en dos
    pasadas; enumerate via walk lineal.
    \item Graph (capítulo~\ref{cap:yannakakis-tree}): bags
    indexados por step (con índices \texttt{by\_origin} y
    \texttt{by\_dest}); join tree construido con GYO; tree-walk
    enumerator.
\end{itemize}

\section{Cómo lo construyen otros transports}

La gracia del AST: cualquier transport que quiera ejecutar un hunt
puede construir una \texttt{Hypothesis} sin pasar por la DSL.
Ejemplos:

\subsection{HTTP REST}

\begin{lstlisting}
POST /run_hunt
{
    "name": "Sysmon W0",
    "dsl": "User -[Auth]-> Host -[Connect]-> IP WITHIN 30m"
}
\end{lstlisting}

\noindent
El handler HTTP toma el campo \texttt{dsl}, llama \texttt{parse\_dsl},
obtiene una \texttt{Hypothesis}, la ejecuta. Como atajo, el
endpoint también acepta una \texttt{Hypothesis} pre-construida en
JSON:

\begin{lstlisting}
POST /run_hunt
{
    "hypothesis": {
        "name": "Sysmon W0",
        "steps": [...],
        "shape": "Chain",
        "window_seconds": 1800,
        "k_simplicity": 1,
        "aggregations": []
    }
}
\end{lstlisting}

\noindent
Bypaseo total del parser. Útil para clientes que generan
hipótesis programáticamente (constructor visual, generators ML).

\subsection{MCP}

\begin{lstlisting}
mcp_call cmd_parse_dsl({"dsl": "..."})  -> Hypothesis (json)
mcp_call cmd_run_hunt({"hypothesis": ...})  -> Vec<HuntResult>
\end{lstlisting}

\noindent
Dos llamadas separadas: una para parse (devuelve la
\texttt{Hypothesis} para inspección), otra para run. Esto permite
que el LLM modifique la \texttt{Hypothesis} programáticamente entre
ambas (ajustar predicados, expandir variables, etc.).

\subsection{Tauri (frontend)}

\begin{lstlisting}
cmd_run_hunt(hypothesis_json: String) -> Vec<HuntResult>
\end{lstlisting}

\noindent
El frontend serializa la \texttt{Hypothesis} (construida a partir
del input del usuario, ya parseada) como JSON. El backend Rust
deserializa y ejecuta. La separación frontend/backend pasa por la
\texttt{Hypothesis}, no por el string DSL.

\section{Estabilidad de la \texttt{Hypothesis}}

\begin{advertencia}
La \texttt{Hypothesis} es la \emph{interfaz pública} del motor para
matchers. Cambios breaking en su shape (renombrar campos, agregar
required fields, cambiar enums) rompen consumers. La política:
agregar campos opcionales es OK; remover o renombrar requiere ADR
+ versionado.
\end{advertencia}

\section{Validación}

Antes de ejecutar, el matcher valida la \texttt{Hypothesis}:

\begin{enumerate}
    \item \texttt{steps} no vacía (al menos 1 step).
    \item Bindings consistentes: si \texttt{step[i].dest\_var ==
    step[j].source\_var}, entonces \texttt{step[i].dest\_type ==
    step[j].source\_type}.
    \item \texttt{shape == Chain} $\Rightarrow$ los \texttt{VarId}s
    forman una secuencia lineal sin saltos.
    \item \texttt{shape == Graph} $\Rightarrow$ chequea
    \(\alpha\)-aciclidad antes de pasar a Yannakakis (sino, se
    cae al DFS).
\end{enumerate}

\noindent
Si alguna validación falla, el matcher devuelve un error
estructurado al transport en lugar de empezar el hunt.

\section{Resumen}

\begin{itemize}
    \item \texttt{Hypothesis}: la representación interna del AST,
    consumida por el matcher en lugar del string DSL.
    \item Campos: \texttt{steps}, \texttt{shape}, \texttt{k\_simplicity},
    \texttt{aggregations}, \texttt{window\_seconds}, \texttt{name}.
    \item \texttt{VarId}: índice numérico para bindings; el matcher
    nunca toca strings.
    \item \texttt{shape == Chain | Graph} dispatcha al
    sub-matcher correspondiente (chain o tree-walk).
    \item Los transports pueden construir \texttt{Hypothesis}
    programáticamente sin pasar por la DSL.
\end{itemize}

\begin{ejercicios}
\begin{enumerate}
    \item \starred{} Construí la \texttt{Hypothesis} (struct Rust)
    para
    \texttt{User c -[Auth]-> Host h1; c -[Auth]-> Host h2}.
    Identificá los \texttt{VarId}s asignados.

    \item ¿Qué pasa si construís una \texttt{Hypothesis} con
    \texttt{shape == Chain} pero los steps NO forman cadena
    (binding shared en el medio)? ¿El matcher lo detecta?

    \item \starred{} Diseñá una API REST para construir
    hipótesis: ¿qué endpoint? ¿Qué payload? ¿Cómo manejás errores
    de validación?

    \item Si la \texttt{Hypothesis} fuera serializable como
    Protobuf en vez de JSON, ¿qué cambiaría? ¿Cuándo conviene?

    \item ¿Por qué los nombres textuales de bindings (\texttt{u},
    \texttt{c}) NO se preservan en \texttt{VarId}? ¿Cuándo importa
    el nombre original?
\end{enumerate}
\end{ejercicios}
